\chapter{Eigenvalues, Eigenvectors, Diagonalization \& SVD}

Eigenvalues and Eigenvectors form the ultimate cornerstone of modern linear algebra and data science. In this capstone chapter, we master \textbf{Eigen-decomposition}, \textbf{Matrix Diagonalization}, and \textbf{Singular Value Decomposition (SVD)}.

\section{Eigenvalues and Eigenvectors}

When a matrix $A \in \mathbb{R}^{n \times n}$ acts on a vector $\mathbf{v}$, it typically changes both the magnitude and direction of $\mathbf{v}$. However, certain non-zero vectors maintain their original direction and are merely scaled by a factor $\lambda$.

\begin{definitionbox}{Eigenvalue Problem}
A non-zero vector $\mathbf{v} \neq \mathbf{0}$ is an \textbf{Eigenvector} of $A$ with corresponding \textbf{Eigenvalue} $\lambda \in \mathbb{R}$ (or $\mathbb{C}$) if:
$$A \mathbf{v} = \lambda \mathbf{v}$$
Equivalently:
$$(A - \lambda I) \mathbf{v} = \mathbf{0}$$
\end{definitionbox}

\begin{theorembox}{The Characteristic Equation}
Since $\mathbf{v} \neq \mathbf{0}$, the matrix $(A - \lambda I)$ must be singular ($\det(A - \lambda I) = 0$).
1. \textbf{Characteristic Polynomial:} $p(\lambda) = \det(A - \lambda I) = 0$. Solve for roots $\lambda_1, \lambda_2, \dots, \lambda_n$.
2. \textbf{Eigenspace:} For each $\lambda_i$, solve $(A - \lambda_i I)\mathbf{v} = \mathbf{0}$ to find the basis of eigenvectors $\text{Null}(A - \lambda_i I)$.
\end{theorembox}

\begin{videocard}{3Blue1Brown: Essence of Linear Algebra (Eigenvalues \& Eigenvectors)}{https://www.youtube.com/watch?v=PFDu9oVAE-g}
Visual intuition showing eigenvectors as axes of rotation/scaling during linear transformations.
\end{videocard}

\section{Matrix Diagonalization}

\begin{definitionbox}{Diagonalization Theorem}
An $n \times n$ matrix $A$ is \textbf{Diagonalizable} if it has $n$ linearly independent eigenvectors $\mathbf{v}_1, \dots, \mathbf{v}_n$.
Form eigenvector matrix $P = [\mathbf{v}_1 \; \dots \; \mathbf{v}_n]$ and diagonal matrix $D = \operatorname{diag}(\lambda_1, \dots, \lambda_n)$. Then:
$$A = P D P^{-1} \iff D = P^{-1} A P$$
\textbf{Matrix Powers:} $A^k = P D^k P^{-1}$, making high-power matrix computations $\mathcal{O}(n)$ instead of $\mathcal{O}(n^k)$!
\end{definitionbox}

\begin{theorembox}{Spectral Theorem for Symmetric Matrices}
If $A = A^T$ is a real symmetric matrix, then:
1. All eigenvalues of $A$ are real numbers.
2. Eigenvectors corresponding to distinct eigenvalues are \textbf{mutually orthogonal}.
3. $A$ is \textbf{orthogonally diagonalizable}: $A = Q D Q^T$, where $Q$ is an orthogonal matrix ($Q^T Q = I$).
\end{theorembox}

\section{Singular Value Decomposition (SVD)}

While Eigen-decomposition requires square matrices $A \in \mathbb{R}^{n \times n}$, \textbf{Singular Value Decomposition (SVD)} applies to \textbf{ANY rectangular matrix} $A \in \mathbb{R}^{m \times n}$.

\begin{formulabox}{Singular Value Decomposition ($A = U \Sigma V^T$)}
Every matrix $A \in \mathbb{R}^{m \times n}$ can be factored into:
$$A = U \Sigma V^T$$
where:
\begin{itemize}
    \item $U \in \mathbb{R}^{m \times m}$: Left singular vectors (orthonormal eigenvectors of $A A^T$).
    \item $\Sigma \in \mathbb{R}^{m \times n}$: Diagonal matrix of singular values $\sigma_1 \ge \sigma_2 \ge \dots \ge \sigma_r > 0$ ($\sigma_i = \sqrt{\lambda_i(A^T A)}$).
    \item $V \in \mathbb{R}^{n \times n}$: Right singular vectors (orthonormal eigenvectors of $A^T A$).
\end{itemize}
\end{formulabox}

\index{Eigenvalue}\index{Eigenvector}\index{SVD}\index{PCA}

\begin{lstlisting}[caption={NumPy Implementation of Eigen-decomposition and SVD}]
import numpy as np

# 1. Eigenvalue Decomposition
A = np.array([[4.0, 1.0], 
              [2.0, 3.0]])

eigenvalues, eigenvectors = np.linalg.eig(A)
print("Eigenvalues:", eigenvalues)
print("Eigenvectors:\n", eigenvectors)

# Verify A @ v = lambda * v for the first eigenvector
v0 = eigenvectors[:, 0]
lam0 = eigenvalues[0]
print("A @ v0:    ", A @ v0)
print("lambda * v0:", lam0 * v0)

# 2. Singular Value Decomposition (SVD)
X = np.array([[3.0, 1.0], 
              [1.0, 3.0], 
              [0.0, 2.0]])

U, S, Vt = np.linalg.svd(X)
print("\nLeft Singular Vectors U shape:", U.shape)
print("Singular Values S:", S)
print("Right Singular Vectors V^T shape:", Vt.shape)
\end{lstlisting}

\begin{videocard}{StatQuest: Singular Value Decomposition (SVD) \& PCA}{https://www.youtube.com/watch?v=gXbThCXj4cA}
Step-by-step visual explanation of SVD matrix factorization and dimensionality reduction.
\end{videocard}

\begin{videocard}{IITM BS Lecture: Mathematics for Data Science II Review}{https://www.youtube.com/watch?v=_hLYDPQtDKY}
Official comprehensive course summary reviewing Linear Algebra and Multivariable Calculus.
\end{videocard}


\newpage
\section{Practice Exercises}
\textit{Detailed step-by-step solutions are available in the Appendix.}

\begin{enumerate}
\item \textbf{Eigenvalues & Eigenvectors ($2 \times 2$):} Find the eigenvalues $\lambda_1, \lambda_2$ and corresponding eigenvectors $\mathbf{v}_1, \mathbf{v}_2$ for the matrix:
    $$A = \begin{bmatrix} 4 & 1 \\ 2 & 3 \end{bmatrix}$$

\item \textbf{Matrix Diagonalization ($A = P D P^{-1}$):} Using the eigenvalues and eigenvectors from Exercise 1, construct the matrices $P$ and $D$ such that $A = P D P^{-1}$. Compute $A^4$ using $P D^4 P^{-1}$.

\item \textbf{Symmetric Matrix Spectral Decomposition:} Verify that $A = \begin{bmatrix} 3 & 1 \\ 1 & 3 \end{bmatrix}$ is symmetric. Find its real eigenvalues and construct an orthogonal matrix $Q$ ($Q^T Q = I$) such that $A = Q D Q^T$.

\item \textbf{Singular Values Calculation:} Given matrix $A = \begin{bmatrix} 3 & 0 \\ 0 & -2 \end{bmatrix}$, compute the matrix $A^T A$, find its eigenvalues $\lambda_i$, and determine the singular values $\sigma_1, \sigma_2$ of $A$.

\item \textbf{Comprehensive Course Review Question (PCA via SVD):} A data matrix $X \in \mathbb{R}^{N \times d}$ has covariance matrix $C = \frac{1}{N} X^T X$.
    \begin{enumerate}
        \item Explain why the principal component directions are the eigenvectors of $C$.
        \item If the eigenvalues of $C$ are $\lambda_1 = 80$, $\lambda_2 = 15$, $\lambda_3 = 5$, what percentage of total dataset variance is explained by keeping the first 2 principal components?
    \end{enumerate}
\end{enumerate}